% !PRINTED none

\begin{center}
\gesperrt{Erster Abschnitt.}

\vspace{1em}
\gesperrt{Allgemeine Gruppentheorie.}

\rule{2cm}{0.4pt}

\vspace{1em}
\S.\,1.

\gesperrt{Definition der Gruppen.}
\end{center}

\noindent Wir haben im ersten Bande bei den Permutationen den Begriff einer Gruppe kennen gelernt und wichtige algebraische Anwendungen von ihm gemacht. Es muss nun unsere n\"achste Aufgabe sein, diesen in der ganzen neueren Mathematik so \"uberaus wichtigen Begriff allgemeiner zu fassen und die dabei herrschenden Gesetze kennen zu lernen. Wir stellen folgende Definition an die Spitze:

Ein System $P$ von Dingen (Elementen) irgend welcher Art wird zur Gruppe, wenn folgende Voraussetzungen erf\"ullt sind:

\noindent 1.\ \gesperrt{Es ist eine Vorschrift gegeben, nach der aus einem ersten und einem zweiten Elemente des Systems ein ganz bestimmtes drittes Element desselben Systems abgeleitet wird.}

Man schreibt symbolisch, wenn $a$ das erste, $b$ das zweite, $c$ das dritte Element ist:
\[
ab = c, \quad c = ab,
\]
und nennt $c$ aus $a$ und $b$ \gesperrt{componirt} und $a$ und $b$ die \gesperrt{Componenten} von $c$.

Bei dieser Composition wird im Allgemeinen nicht das commutative Gesetz vorausgesetzt, d.\,h. es kann $ab$ von $ba$ verschieden sein, dagegen wird

\noindent 2.\ \gesperrt{das associative Gesetz vorausgesetzt,} d.\,h. wenn $a, b, c$ irgend drei Elemente aus $P$ sind, so ist
\[
(ab)c = a(bc),
\]
%CONT>>
